% ---- pg-161 （书页 149） ----

对许多 $\lambda$ 值计算 $l(\lambda)$, 对计算机来讲是件容易的工作. 为了研究接近最大值的 $\hat{\lambda}$ 的特性, 作出 $l(\lambda)$ 与 $\lambda$ 的图形, 并把数对 $(\lambda, l(\lambda))$ 表示成表格, 是有用的. 例如, 如果 $\lambda = 0$(对数)或 $\lambda = 1/2$(平方根)靠近 $\hat{\lambda}$, 由于其简单性, 其中之一会被优先选用.

有些统计学家建议, 与其进行\hyperref[eq:4-35]{式(4-35)}中的计算, 不如用一种等价的方法, 即固定 $\lambda$ 值, 建立新变量
\phantomsection\label{eq:4-37}
\begin{equation}
y_j^{(\lambda)} = \frac{x_j^\lambda - 1}{\lambda\left[\left(\prod_{i=1}^{n} x_i\right)^{1/n}\right]^{\lambda-1}} \qquad j = 1, \cdots, n \tag{4-37}
\end{equation}
然后计算样本方差. 方差最小时的 $\lambda$ 值与使\hyperref[eq:4-35]{式(4-35)}达到最大时的 $\lambda$ 是相同的.

\mvbold{评注}\quad 现在已知道, 通过使 $l(\lambda)$ 最大化而得到的变换通常会改进正态性近似. 但是, 连 $\lambda$ 的最好的选择都不能保证将产生一个充分符合正态分布的变换数据集. 按照\hyperref[eq:4-35]{式(4-35)}选择的变换所产生的结果总应当仔细地检查, 因为有可能违背尝试性的正态性假定. 这个警告适用于根据任何其他方法选择的变换.

\noindent\textbf{例 4.16}\quad (对一元数据确定幂变换)

在例 4.10 中, 我们给出了 $n = 42$ 个微波炉在炉门关闭时的辐射量读数. 在\hyperref[fig:4.6]{图 4.6} 中, 这些数据的 Q-Q 图表明观测值背离了如果它们是正态分布时的期望情形. 由于所有的观测值为正, 我们对数据进行幂变换, 希望它会产生更接近正态的结果. 把我们的注意力限制在\hyperref[eq:4-34]{式(4-34)}中给出的一类变换, 我们必须寻求 $\lambda$ 的值, 使\hyperref[eq:4-35]{式(4-35)}中的函数 $l(\lambda)$ 达到最大值.

对于 $\lambda$ 的若干取值, 数对 $(\lambda, l(\lambda))$ 列于下表:

\begin{mvtable}
\begin{tabular}{cc|cc}
\toprule
$\lambda$ & $l(\lambda)$ & $\lambda$ & $l(\lambda)$ \\
\midrule
$-1.00$ & 70.52 & 0.40 & 106.20 \\
$-0.90$ & 75.65 & 0.50 & 105.50 \\
$-0.80$ & 80.46 & 0.60 & 104.43 \\
$-0.70$ & 84.94 & 0.70 & 103.03 \\
$-0.60$ & 89.06 & 0.80 & 101.33 \\
$-0.50$ & 92.79 & 0.90 & 99.34 \\
$-0.40$ & 96.10 & 1.00 & 97.10 \\
$-0.30$ & 98.97 & 1.10 & 94.64 \\
$-0.20$ & 101.39 & 1.20 & 91.96 \\
$-0.10$ & 103.35 & 1.30 & 89.10 \\
$0.00$ & 104.83 & 1.40 & 86.07 \\
$0.10$ & 105.84 & 1.50 & 82.88 \\
$0.20$ & 106.39 & & \\
\fbox{0.30} & \fbox{106.51} & & \\
\bottomrule
\end{tabular}
\end{mvtable}
